\documentclass[landscape,twocolumn,letter]{book}
%\documentclass[10pt,letter]{scrartcl}
\usepackage[spanish]{babel}
\usepackage[left=1.5cm, right=1.5cm, top=1.5cm, bottom=1.5cm]{geometry}

%\usepackage{geometry}
%\usepackage{amsmath}
%\usepackage{mathrsfs}
%\usepackage{graphicx}
%\usepackage{caption}
%\usepackage{subcaption}
%\usepackage[utf8]{inputenc}
\usepackage{hyperref}

\usepackage{graphicx}
\usepackage{caption}
\usepackage{subcaption}
\usepackage{setspace}
\usepackage{fancyhdr}
\usepackage{mathtools}

\usepackage{setspace}
%\spacing{1.5}

\DeclareCaptionLabelSeparator{emdash}{\textemdash}
\captionsetup[figure]{labelsep=emdash,font=onehalfspacing,position=bottom}
\captionsetup{%
	singlelinecheck=off,
	skip=2pt,
	justification=centering,
}

%\usepackage[T1]{fontenc}
%\usepackage{imakeidx}
\makeindex 

\title{Apuntes del Instructor: Electrónica III (PROBLEMARIO)}
\date{2026 Mayo\\rev 0.1k}
\author{Alberto Medrano\thanks{Instructor, Universidad de Los Andes (ULA)}, Manuel Rodriguez\thanks{Prof. Titular, Universidad de Los Andes (ULA)}, Roman Alí Contreras\thanks{Prof. Titular, Universidad de Los Andes (ULA)}}
\begin{document}
	\pagestyle{fancy}
	
	\fancyhead{} % clear all header fields
	\fancyhead[RO,LE]{\textbf{Apuntes del Instructor:} Electrónica III}
	\fancyfoot{} % clear all footer fields
	\fancyfoot[LE,RO]{\thepage}
	\fancyfoot[LO,CE]{Compilado por: A. Medrano}
	
	\maketitle
	
	\includegraphics[width=0.4\linewidth]{_fig/ULA.pdf}
	
	\textbf{Sobre el contenido de éstos apuntes}
	
	- - -
	
	Éstos son apuntes personales para la materia Electrónica III, contenido influenciado del Libro ELECTRÓNICA LINEAL AVANZADA, Asignatura: Electrónica III, Autor: Manuel Rodríguez y con contribuciones de ejercicios del Prof. Román Ali compilados en ésta guía. También considera algún material didáctico para complementar la información aquí proporcionada. De encontrar algún error o sugerencia por favor notificar vía correo electrónico a \href{mailto: medrano.alberto@ula.ve}{medrano.alberto@ula.ve}
	
	\tableofcontents
	
	\chapter{La respuesta de frecuencia de circuitos electrónicos} %p2
	
	\chapter{Amplificadores de varias etapas} %p70
	
	\chapter{Retroalimentación} %p129
	
	\chapter{Estabilidad y osciladores} %p183
	
	\chapter{Lazo de enganche de fase (PLL)} %p223
	
	\chapter{Amplificadores de potencia} %p241

	
\end{document}
